\chapter{Qubitization and Quantum Signal Processing}
% Solutions to the exercises of \cref{ch:lec08}.

\begin{solution}[Exercise~\ref{ex:lec08-1}]
With $\lambda=\cos\theta$ and $\sqrt{1-\lambda^2}=\sin\theta\ (\ge0$ on $[0,\pi])$,
\[
  O(\lambda)=U(\lambda)Z
  =\begin{pmatrix}\lambda&\sqrt{1-\lambda^2}\\ \sqrt{1-\lambda^2}&-\lambda\end{pmatrix}
   \begin{pmatrix}1&0\\0&-1\end{pmatrix}
  =\begin{pmatrix}\lambda&-\sqrt{1-\lambda^2}\\ \sqrt{1-\lambda^2}&\lambda\end{pmatrix}
  =\begin{pmatrix}\cos\theta&-\sin\theta\\ \sin\theta&\cos\theta\end{pmatrix},
\]
a rotation by $\theta$; hence $O(\lambda)^k$ is the rotation by $k\theta$, with
$(1,1)$-entry $\cos(k\theta)=T_k(\lambda)$ and $(2,1)$-entry
$\sin(k\theta)=\sin\theta\,U_{k-1}(\lambda)$. From
$\cos((k+1)\theta)+\cos((k-1)\theta)=2\cos\theta\cos(k\theta)$ we get
$T_{k+1}(\lambda)+T_{k-1}(\lambda)=2\lambda T_k(\lambda)$, i.e.\ the recurrence.
Cubing,
\[
  O(\lambda)^3=\begin{pmatrix}\cos3\theta&-\sin3\theta\\ \sin3\theta&\cos3\theta\end{pmatrix},
  \quad
  T_3(\lambda)=\cos3\theta=4\lambda^3-3\lambda,\quad
  U_2(\lambda)=\frac{\sin3\theta}{\sin\theta}=4\lambda^2-1,
\]
using $\cos3\theta=4\cos^3\theta-3\cos\theta$ and
$\sin3\theta=(3-4\sin^2\theta)\sin\theta=(4\cos^2\theta-1)\sin\theta$.
\end{solution}

\begin{solution}[Exercise~\ref{ex:lec08-2}]
Here $A=\tfrac12\Xgate=\tfrac12(\proj{+}-\proj{-})$ has eigenvalues
$\lambda_\pm=\pm\tfrac12$ with eigenvectors $\ket\pm=\tfrac1{\sqrt2}(\ket0\pm\ket1)$,
and $m=1$ so $Z_\Pi=\Zgate$ on the single flag qubit. By \cref{lem:lec08-qubit},
\[
  U_A\ket0\ket\pm=\pm\tfrac12\ket0\ket\pm+\tfrac{\sqrt3}{2}\ket{\perp_\pm},
  \qquad
  [U_A]_{\mathcal B_\pm}
  =\begin{pmatrix}\pm\tfrac12&\tfrac{\sqrt3}{2}\\[2pt]\tfrac{\sqrt3}{2}&\mp\tfrac12\end{pmatrix},
\]
where $\sqrt{1-\lambda_\pm^2}=\sqrt{1-\tfrac14}=\tfrac{\sqrt3}{2}$. Then
$[U_A\Zgate]_{\mathcal B_\pm}=\bigl(\begin{smallmatrix}\pm\frac12&-\frac{\sqrt3}2\\
\frac{\sqrt3}2&\pm\frac12\end{smallmatrix}\bigr)=O(\lambda_\pm)$, whose $(1,1)$-entry
is $\lambda_\pm=T_1(\lambda_\pm)$; summing over the two eigenspaces,
$\Pi(U_A\Zgate)\Pi=\proj0\otimes\sum_\pm\lambda_\pm\proj\pm=\proj0\otimes A$, so
$k=1$ gives $T_1(A)=A$. For $k=2$, $[O(\lambda_\pm)^2]_{11}=2\lambda_\pm^2-1
=2\cdot\tfrac14-1=-\tfrac12$ for \emph{both} signs, so
$\Pi(U_A\Zgate)^2\Pi=\proj0\otimes(-\tfrac12)(\proj++\proj-)
=\proj0\otimes(-\tfrac12\id)=\proj0\otimes T_2(A)$, matching
$T_2(A)=2A^2-\id=-\tfrac12\id$ of \cref{exa:lec08-T2}.
\end{solution}

\begin{solution}[Exercise~\ref{ex:lec08-3}]
The matrix $M\defeq[U_A]_{\mathcal B_j}=\bigl(\begin{smallmatrix}\lambda&\sqrt{1-\lambda^2}
\\ \sqrt{1-\lambda^2}&-\lambda\end{smallmatrix}\bigr)$ satisfies $M=M\dg$ and
$M^2=(\lambda^2+(1-\lambda^2))\id=\id$, with $\det M=-\lambda^2-(1-\lambda^2)=-1$: it
is a reflection. Right-multiplying by $Z=[Z_\Pi]_{\mathcal B_j}$ gives
$MZ=O(\lambda_j)$ with $\det(MZ)=(\det M)(\det Z)=(-1)(-1)=+1$ and $(MZ)(MZ)\dg=\id$:
a rotation. Powers of $U_A$ alone give $[U_A^k]_{\mathcal B_j}=M^k$, which cycles
$\id,M,\id,M,\dots$ and never produces $\cos(k\theta)=T_k(\lambda)$ for $k\ge2$ (a
reflection has eigenvalues $\pm1$, so its powers only ever give $\pm1$ on the block,
not the continuum of Chebyshev values). Interposing the reflection $Z_\Pi$ converts
each reflection $M$ into the rotation $MZ$, and it is the \emph{minimal} fix because
a single extra determinant-$(-1)$ factor per query is exactly what is needed to make
the product a rotation ($\det=+1$); using no reflection leaves a reflection, and
using two per query would return a reflection again.
\end{solution}

\begin{solution}[Exercise~\ref{ex:lec08-4}]
By \cref{thm:lec08-general}, even Chebyshev polynomials come from $k$ pairs
$U_A\dg Z_\Pi\,U_A Z_\Pi$. For $T_4(A)$ take $k=2$:
\[
  T_4:\quad \bigl(U_A\dg Z_\Pi\,U_A Z_\Pi\bigr)^2
   =U_A\dg Z_\Pi\,U_A Z_\Pi\,U_A\dg Z_\Pi\,U_A Z_\Pi,
\]
which uses $2$ queries to $U_A$ and $2$ to $U_A\dg$, total $4$ signal queries
$=\deg T_4$. For the odd $T_3(A)$ prepend one $U_A Z_\Pi$ to the $k=1$ pair:
\[
  T_3:\quad U_A Z_\Pi\,\bigl(U_A\dg Z_\Pi\,U_A Z_\Pi\bigr)
   =U_A Z_\Pi\,U_A\dg Z_\Pi\,U_A Z_\Pi,
\]
which uses $2$ queries to $U_A$ and $1$ to $U_A\dg$, total $3=\deg T_3$. In both
cases the number of signal operators (queries to $U_A$ or $U_A\dg$) equals the
Chebyshev degree; the $U_A,U_A\dg$ alternation is what lets each query advance the
degree by one while respecting the map $\mathcal H_j\leftrightarrow\mathcal H_j'$.
\end{solution}

\begin{solution}[Exercise~\ref{ex:lec08-5}]
Compute
\[
  U_\Phi(\lambda)=\ee^{\ii\phi_0 Z}O(\lambda)\ee^{\ii\phi_1 Z}
  =\begin{pmatrix}\ee^{\ii\phi_0}&0\\0&\ee^{-\ii\phi_0}\end{pmatrix}
   \begin{pmatrix}\lambda&-\sqrt{1-\lambda^2}\\ \sqrt{1-\lambda^2}&\lambda\end{pmatrix}
   \begin{pmatrix}\ee^{\ii\phi_1}&0\\0&\ee^{-\ii\phi_1}\end{pmatrix}
  =\begin{pmatrix}\ee^{\ii(\phi_0+\phi_1)}\lambda & -\ee^{\ii(\phi_0-\phi_1)}\sqrt{1-\lambda^2}\\[3pt]
    \ee^{-\ii(\phi_0-\phi_1)}\sqrt{1-\lambda^2}&\ee^{-\ii(\phi_0+\phi_1)}\lambda\end{pmatrix}.
\]
Comparing with \eqref{eq:lec08-qspform}, $P(\lambda)=\ee^{\ii(\phi_0+\phi_1)}\lambda$
and $Q(\lambda)=\ee^{\ii(\phi_1-\phi_0)}$ (constant, so that
$-Q\sqrt{1-\lambda^2}$ matches the top-right entry). Conditions:
$\deg P=1\le1$, $\deg Q=0\le0$ (1); $P$ odd, $Q$ even (2), with $d=1$; and
$\abs{P}^2+(1-\lambda^2)\abs{Q}^2=\lambda^2+(1-\lambda^2)=1$ (3). Taking
$\phi_0=\phi_1=0$ gives $P(\lambda)=\lambda$; taking
$\phi_0=\phi_1=\tfrac\pi4$ gives $P(\lambda)=\ee^{\ii\pi/2}\lambda=\ii\lambda$. A
degree-$1$ sequence \emph{cannot} give $P(\lambda)=\tfrac12\lambda$: condition (3)
forces $\abs{P(1)}=1$, whereas $\tfrac12\lambda$ has $\abs{P(1)}=\tfrac12$. This does
not contradict \cref{cor:lec08-real}: there one asks for $\operatorname{Re}P=p$ with
$\abs{p}\le1$, and the subnormalized target $p(\lambda)=\tfrac12\lambda$ is realized
as the \emph{real part} of a valid $P$ (of the same degree, with $\abs{P}\le1$),
e.g.\ via the $\tfrac12(P+P^*)$ LCU construction, not as $P$ itself.
\end{solution}

\begin{solution}[Exercise~\ref{ex:lec08-6}]
Multiply the degree-$(d-1)$ form by $O(\lambda)\ee^{\ii\phi_d Z}$. First,
\[
  \begin{pmatrix}P&-Q\sqrt{1-\lambda^2}\\ Q^*\sqrt{1-\lambda^2}&P^*\end{pmatrix}
  \begin{pmatrix}\lambda&-\sqrt{1-\lambda^2}\\ \sqrt{1-\lambda^2}&\lambda\end{pmatrix}
  =\begin{pmatrix}\lambda P-(1-\lambda^2)Q & -\sqrt{1-\lambda^2}(P+\lambda Q)\\[3pt]
    \sqrt{1-\lambda^2}(P^*+\lambda Q^*)&\lambda P^*-(1-\lambda^2)Q^*\end{pmatrix},
\]
using $-Q\sqrt{1-\lambda^2}\cdot\sqrt{1-\lambda^2}=-(1-\lambda^2)Q$ in the $(1,1)$
slot and $P\cdot(-\sqrt{1-\lambda^2})+(-Q\sqrt{1-\lambda^2})\lambda
=-\sqrt{1-\lambda^2}(P+\lambda Q)$ in the $(1,2)$ slot. Right-multiplying by
$\ee^{\ii\phi_d Z}=\mathrm{diag}(\ee^{\ii\phi_d},\ee^{-\ii\phi_d})$ scales column $1$
by $\ee^{\ii\phi_d}$ and column $2$ by $\ee^{-\ii\phi_d}$, giving
$\widetilde P=\ee^{\ii\phi_d}(\lambda P-(1-\lambda^2)Q)$ and (from the top-right
entry $-\widetilde Q\sqrt{1-\lambda^2}=-\ee^{-\ii\phi_d}\sqrt{1-\lambda^2}(P+\lambda
Q)$) $\widetilde Q=\ee^{-\ii\phi_d}(P+\lambda Q)$, which is \eqref{eq:lec08-newPQ}.
\emph{Degrees:} $\deg\widetilde P\le\max(1+\deg P,\,2+\deg Q)\le\max(d,\,2+(d-2))=d$
and $\deg\widetilde Q\le\max(\deg P,\,1+\deg Q)\le d-1$. \emph{Parity:} $P$ has parity
$(d-1)$ and $Q$ parity $(d-2)$; multiplying by $\lambda$ (odd) or $(1-\lambda^2)$
(even) makes both terms of $\widetilde P$ have parity $d$, and both terms of
$\widetilde Q$ parity $d-1$. \emph{Unitarity:} the map
$(P,Q)\mapsto(\lambda P-(1-\lambda^2)Q,\,P+\lambda Q)$ followed by unit-modulus phases
preserves $\abs{\cdot}^2$ because $O(\lambda)$ and $\ee^{\ii\phi_d Z}$ are unitary;
explicitly
\[
  \abs{\widetilde P}^2+(1-\lambda^2)\abs{\widetilde Q}^2
  =\abs{\lambda P-(1-\lambda^2)Q}^2+(1-\lambda^2)\abs{P+\lambda Q}^2
  =(\lambda^2+(1-\lambda^2))\bigl(\abs P^2+(1-\lambda^2)\abs Q^2\bigr),
\]
where the cross terms $-2\lambda(1-\lambda^2)\operatorname{Re}(P\bar Q)$ and
$+2\lambda(1-\lambda^2)\operatorname{Re}(P\bar Q)$ cancel; the bracket is $1$ by the
inductive hypothesis, so condition (3) is preserved.
\end{solution}

\begin{solution}[Exercise~\ref{ex:lec08-7}]
With $\Phi=(0,0,0)$, $U_\Phi(\lambda)=O(\lambda)^2$, and by \cref{lem:lec08-rot} its
$(1,1)$-entry is $\cos2\theta=2\cos^2\theta-1=2\lambda^2-1=T_2(\lambda)$; explicitly
\[
  O(\lambda)^2=\begin{pmatrix}\lambda&-\sqrt{1-\lambda^2}\\ \sqrt{1-\lambda^2}&\lambda\end{pmatrix}^2
  =\begin{pmatrix}2\lambda^2-1&-2\lambda\sqrt{1-\lambda^2}\\ 2\lambda\sqrt{1-\lambda^2}&2\lambda^2-1\end{pmatrix}.
\]
For the eigenvalue-transform convention, recall $O(\lambda)=U_A(\lambda)Z$ and
$\ee^{\ii\frac\pi2 Z}=\ii Z$, so $O(\lambda)\ee^{\ii\phi Z}=U_A(\lambda)Z\ee^{\ii\phi Z}
=(-\ii)\,U_A(\lambda)\,\ee^{\ii(\phi+\frac\pi2)Z}$ after writing
$Z=(-\ii)\ee^{\ii\frac\pi2 Z}$ and combining exponents. Applying this to each of the
$d=2$ signal factors of $O(\lambda)^2$ (phases $0,0$) converts the all-zero
$O$-convention into the $e$-convention with
$e=(0+\tfrac\pi4,\,0+\tfrac\pi2,\,0+\tfrac\pi4)=(\tfrac\pi4,\tfrac\pi2,\tfrac\pi4)$
(the interior phases pick up $\tfrac\pi2$ and the two boundary phases $\tfrac\pi4$
each, the standard QET convention), and pulls out the global factor $(-\ii)^2=-1$. Concretely one checks by
direct multiplication that
\[
  \ee^{\ii\frac\pi4 Z}U_A(\lambda)\ee^{\ii\frac\pi2 Z}U_A(\lambda)\ee^{\ii\frac\pi4 Z}
  =(-\ii)^2\begin{pmatrix}2\lambda^2-1&*\\ *&*\end{pmatrix}
  =-\begin{pmatrix}T_2(\lambda)&*\\ *&*\end{pmatrix},
\]
so up to the overall phase $(-\ii)^d=(-\ii)^2=-1$ the top-left block is
$T_2(\lambda)=2\lambda^2-1$, as claimed in \cref{exa:lec08-qsp}(c). (The phase
$(-\ii)^d$ is a physically irrelevant global factor and is dropped in the block
encoding.)
\end{solution}
